\section{Trigonometry}

All angles are measured in radians unless stated otherwise. Euler's formula provides the useful bridge between trigonometric and exponential functions:
\begin{equation}
e^{\mathrm{i}\theta}=\cos\theta+\mathrm{i}\sin\theta,
\end{equation}
so that
\begin{equation}
\cos\theta=\frac{e^{\mathrm{i}\theta}+e^{-\mathrm{i}\theta}}{2},
\qquad
\sin\theta=\frac{e^{\mathrm{i}\theta}-e^{-\mathrm{i}\theta}}{2\mathrm{i}}.
\end{equation}

The basic identities are
\begin{equation}
\sin^2\theta+\cos^2\theta=1,
\qquad
1+\tan^2\theta=\sec^2\theta,
\qquad
1+\cot^2\theta=\csc^2\theta.
\end{equation}
The addition and double-angle formulae are
\begin{align}
\sin(\alpha\pm\beta)&=\sin\alpha\cos\beta\pm\cos\alpha\sin\beta,\\
\cos(\alpha\pm\beta)&=\cos\alpha\cos\beta\mp\sin\alpha\sin\beta,\\
\tan(\alpha\pm\beta)&=\frac{\tan\alpha\pm\tan\beta}{1\mp\tan\alpha\tan\beta},\\
\sin 2\alpha&=2\sin\alpha\cos\alpha,\\
\cos 2\alpha&=\cos^2\alpha-\sin^2\alpha=2\cos^2\alpha-1=1-2\sin^2\alpha,\\
\tan 2\alpha&=\frac{2\tan\alpha}{1-\tan^2\alpha}.
\end{align}
The half-angle formulae read
\begin{equation}
\sin^2\frac{\alpha}{2}=\frac{1-\cos\alpha}{2},
\qquad
\cos^2\frac{\alpha}{2}=\frac{1+\cos\alpha}{2},
\qquad
\tan\frac{\alpha}{2}=\frac{\sin\alpha}{1+\cos\alpha}=\frac{1-\cos\alpha}{\sin\alpha}.
\end{equation}

Product-to-sum and sum-to-product formulae are especially useful in Fourier calculations:
\begin{align}
\sin\alpha\cos\beta&=\frac12\left[\sin(\alpha+\beta)+\sin(\alpha-\beta)\right],\\
\cos\alpha\sin\beta&=\frac12\left[\sin(\alpha+\beta)-\sin(\alpha-\beta)\right],\\
\cos\alpha\cos\beta&=\frac12\left[\cos(\alpha+\beta)+\cos(\alpha-\beta)\right],\\
\sin\alpha\sin\beta&=-\frac12\left[\cos(\alpha+\beta)-\cos(\alpha-\beta)\right],\\
\sin\alpha+\sin\beta&=2\sin\frac{\alpha+\beta}{2}\cos\frac{\alpha-\beta}{2},\\
\cos\alpha+\cos\beta&=2\cos\frac{\alpha+\beta}{2}\cos\frac{\alpha-\beta}{2}.
\end{align}
Finally,
\begin{equation}
\arcsin x+\arccos x=\frac{\pi}{2},
\qquad
\arctan x+\operatorname{arccot}x=\frac{\pi}{2}\quad(x\in\mathbb{R}).
\end{equation}
For contour integrals involving rational functions of sine and cosine, setting $z=e^{\mathrm{i}\theta}$ gives $\cos\theta=(z+z^{-1})/2$ and $\sin\theta=(z-z^{-1})/(2\mathrm{i})$.
